% Generated draft: requires mathematical and source review.
\documentclass[11pt]{article}
\usepackage[margin=25mm]{geometry}
\usepackage{fontspec}
\setmainfont{Latin Modern Roman}
\usepackage{amsmath,amssymb,mathtools,unicode-math}
\setmathfont{Latin Modern Math}
\usepackage{enumitem,needspace,xurl,hyperref}
\setlength{\parindent}{0pt}
\setlength{\parskip}{5pt}
\setlength{\emergencystretch}{3em}
\newcommand{\N}{\mathbb{N}}
\newcommand{\Z}{\mathbb{Z}}
\newcommand{\Q}{\mathbb{Q}}
\newcommand{\R}{\mathbb{R}}
\newcommand{\C}{\mathbb{C}}
\newcommand{\F}{\mathbb{F}}
\newcommand{\A}{\mathbb{A}}
\newcommand{\PP}{\mathbb{P}}
\newcommand{\E}{\mathbb{E}}
\newcommand{\eps}{\varepsilon}
\newcommand{\dd}{\,\mathrm{d}}
\begin{document}
\section*{UnsolvedMath ID: 700762}
\textbf{Source problem code:} catalogue entry
\subsection{Definitions and mathematical statement}
The source assumes that \(F\) is an ``infinite field or a skew-field.'' In the formulation below, a skew-field means a division ring, and infinitude is read as applying to \(F\) in both cases. The wording does not explicitly repeat the infinitude qualification for skew-fields.

For the usual positive-dimensional formulation, let \(d\geq 1\) be an integer. The source does not specify the range of \(d\). Use right vector spaces over \(F\), and put \(V=F^{d+1}\). The projective space \(PG(d,F)\) consists of the one-dimensional right \(F\)-subspaces of \(V\), with its projective subspaces induced by right linear subspaces of \(V\). Define
\[
PGL(d+1,F)=GL_{d+1}(F)/\{zI_{d+1}:z\in Z(F)^\times\},
\]
where \(Z(F)\) is the centre of \(F\). This group acts on \(PG(d,F)\) by its natural action on subspaces. A complete flag is a chain
\[
0<V_1<V_2<\cdots<V_d<V,
\qquad \dim_F V_i=i.
\]
A subgroup is flag-transitive if it acts transitively on these complete flags. Here complete flags are the conventional interpretation of the source's unqualified word ``flag.'' A subgroup is \(2\)-transitive on points if, for any points \(p\ne q\) and \(p'\ne q'\), some element sends \(p\) to \(p'\) and \(q\) to \(q'\).

\begin{enumerate}
\item[(a)] Determine all subgroups \(H\leq PGL(2,F)\) that are transitive on the projective line
\[
PG(1,F)\cong F\cup\{\infty\}.
\]
\item[(c)] Determine all flag-transitive subgroups \(H\leq PGL(d+1,F)\).
\item[(d)] Determine all subgroups \(H\leq PGL(d+1,F)\) that are \(2\)-transitive on the points of \(PG(d,F)\).
\end{enumerate}

The following part is retained with its supplied archive qualification:
\begin{enumerate}
\item[(b)] For a subring \(R\subseteq F\), find necessary and sufficient conditions on \(R\) ensuring that the group denoted in the source by \(PGL(d+1,R)\) is flag-transitive on \(PG(d,F)\).
\end{enumerate}
The meaning of this group and its action requires clarification for general subrings of skew-fields. If \(R\) contains the identity of \(F\), a precise matrix-action interpretation is to ask for flag-transitivity of
\[
G_R=\pi_F\bigl(GL_{d+1}(R)\bigr)\leq PGL(d+1,F),
\]
where \(\pi_F:GL_{d+1}(F)\to PGL(d+1,F)\) is the natural quotient map and \(GL_{d+1}(R)\) consists of matrices invertible over \(R\). Under this interpretation, part (b) asks for conditions equivalent to
\[
\forall\,\mathcal F,\mathcal F'\text{ complete flags of }PG(d,F),
\quad \exists g\in G_R:\ g\mathcal F=\mathcal F'.
\]
This is a conditional interpretation, since the source does not specify a common-identity convention for subrings or a definition of \(PGL(d+1,R)\). In particular, scalars central in \(R\) need not be central in \(F\), so quotienting by all central scalar units of \(R\) does not automatically give the stated action on \(PG(d,F)\).

The supplied record labels the overall problem partially solved, places (a), (c), and (d) in the main section, and reports that necessary and sufficient conditions for (b) were given by S. A. Zyubin. These are supplied source-status assertions; no conditions or proof from that work are supplied here.
\subsection{Short English statement}
Determine the transitive projective linear groups on a projective line, the flag-transitive projective linear groups in dimension \(d\), and those that are \(2\)-transitive on points. The archived part asks when a subring produces a flag-transitive group and is reported solved by Zyubin. The dimension range, flag convention, and subring-group notation require clarification, particularly over skew-fields.
\subsection{Sources}
\begin{itemize}
\item Local UnsolvedMath registry, permanent ID \(700762\), exact source problem code \texttt{catalogue entry}, title \emph{Kourovka Notebook Problem 11.70}; supplied upstream revision \texttt{d371ee917ac9647fbb06cd6fc529deec6f3dda82}.
\item P. M. Neumann and C. E. Praeger, proposers of Problem 11.70 in \emph{The Kourovka Notebook (Unsolved Problems in Group Theory)}, 11th issue (1990). The supplied background identifies an update in the 21st issue (2026), \texttt{arXiv:1401.0300v48}. The supplied collection description names E. I. Khukhro and V. D. Mazurov as editors.
\item Supplied Archive of Solved Problems attribution for part (b): S. A. Zyubin, \emph{Siberian Electron. Math. Rep.} \textbf{11} (2014), 64--69 (Russian). The archive reports that this work gives necessary and sufficient conditions. No article title, explicit conditions, or external URL is supplied.
\item The supplied status review was checked on \texttt{2026-10-16}; it labels the problem \texttt{partially\_solved}, identifies part (b) as archived solved, and says the remaining parts are in the main section. This is supplied dated curation, not an independently verified literature assessment.
\end{itemize}
\end{document}
